% ═════════════════════════════════════════════════════════════════════════════
\section{Conclusion and Open Problems}
\label{sec:conclusion}
% ═════════════════════════════════════════════════════════════════════════════

Parameterized verification asks for guarantees that hold for every number of
components, a question that is undecidable in general. The techniques surveyed
in this chapter respond to this difficulty in different ways, but the same few
ideas recur across all of them: a finite symbolic
representation of infinite families of states, monotonicity with respect to a
well-quasi-ordering, and the observation that the behavior of a large instance
is often already visible in a small one. The applications to weak memory
models and to quantum circuits show that these ideas carry over to settings
that look very different from the protocols for which they were first
developed.

Several problems remain open or only partially solved. Liveness properties
under fairness are much less well supported than safety properties, both in
theory and in tools, since abstractions that preserve safety often lose the
information needed for fairness. For weak memory models, parameterized
verification is well understood for TSO and some of its extensions, but the
more relaxed models of ARM and POWER processors, and the language-level models
of C11 and Java, are only partly covered. For distributed protocols, the
automatic inference of invariants with alternating quantifiers, which are
needed for many consensus protocols, remains difficult, and current tools
handle it by restricting the logic or by relying on the user. In the quantum
setting, the automata-based approach covers parameterized circuits, and its
extension to programs with measurement, classical control, and noise is a
natural next step. Finally, the parameterized view applies to further classes of
systems built from many copies of the same component, such as multi-agent
systems and large-scale distributed services, where the same questions about
correctness for any number of components arise.
